\section{K\"ahler models and symplectic transfer}\label{geo:devices}

We next explain how compact packings in projective normal models
transfer to a prescribed target form while avoiding specified divisors.
The final two constructions identify the models used below and the
shell in which the smaller balls will lie when one ball is large.

We use Chern class units: a projective line has area one for the
Fubini--Study form representing \(c_1(\cO_{\PP^n}(1))\).
An ample rational polarization means an element of
\(\operatorname{Pic}(X)\otimes_{\Z}\mathbb Q\) with an ample positive
integral multiple. We use the same notation for its real Chern class,
and omit pullbacks when their meaning is clear.
Projective bundles parametrize lines, and
\[
 U=c_1\bigl(\cO_{\PP(E)}(1)\bigr)
\]
is the dual tautological class. Our normalization of \(\dd^c\) is
\(\dd^c\log|z|^2=\dd\theta\), where
\(\theta=\arg z/(2\pi)\). Thus
\(\dd\dd^c\log|z|^2\) has unit mass at the origin.

\subsection{Moser isotopies with prescribed invariant submanifolds}

We use Moser's deformation method \cite[Section~4]{Moser1965},
with a normal first-jet correction to preserve the specified
complex submanifolds as sets.

\begin{lemma}[Relative K\"ahler Moser lemma]\label{geo:relative-moser}
Let \(X\) be a compact complex manifold and let \(\omega_0,\omega_1\)
be cohomologous K\"ahler forms. Let \(S_1,\ldots,S_e\) be pairwise
disjoint closed smooth complex submanifolds. There is an isotopy
\(\psi_t\) of \(X\), starting at the identity, such that
\[
 \psi_1^*\omega_1=\omega_0,\qquad \psi_t(S_j)=S_j.
\]
If a compact Lie group \(G\) acts holomorphically on \(X\), the forms
are \(G\)-invariant, and every \(S_j\) is \(G\)-invariant, the isotopy
can be chosen \(G\)-equivariant. It then preserves every connected
component of the fixed locus \(X^G\), whether or not that component
meets any \(S_j\).
\end{lemma}

\begin{proof}
Put \(\omega_t=(1-t)\omega_0+t\omega_1\). Each \(\omega_t\) is
K\"ahler. Choose a smooth real one-form \(\alpha\) such that
\(\dd\alpha=\omega_1-\omega_0\).
Since \(S_j\) is complex, its tangent bundle is symplectic for
\(\omega_t\), and
\[
 TX|_{S_j}=TS_j\oplus N_{j,t},\qquad
 N_{j,t}=(TS_j)^{\omega_t}.
\]
These splittings depend smoothly on \(t\). Along \(S_j\), prescribe
the first jet of a function \(f_t\) by requiring
\[
 f_t|_{S_j}=0,\qquad
 \dd f_t|_{TS_j}=0,\qquad
 \dd f_t|_{N_{j,t}}=-\alpha|_{N_{j,t}}.
\]
The direct-sum decomposition makes this a well-defined smooth
covector field. It is realized by a smooth function in a fixed
tubular neighborhood: evaluate the prescribed covector on the
normal variable and multiply by a cutoff equal to one near the
zero section. Compactness of the parameter interval permits fixed
neighborhoods and cutoffs. The neighborhoods for different \(S_j\)
may be chosen disjoint. Adding the resulting functions produces
a global smooth family \(f_t\).

Set \(\alpha_t=\alpha+\dd f_t\), and solve
\[
 \iota_{X_t}\omega_t=-\alpha_t.
\]
Along \(S_j\), the right side annihilates \(N_{j,t}\), so
\(X_t\) is tangent to \(S_j\).
Its flow exists throughout \(0\le t\le1\), because \(X\) is compact,
and the usual Moser calculation gives
\[
 \frac{\dd}{\dd t}(\psi_t^*\omega_t)
 =\psi_t^*\bigl(\dd\alpha+\dd\iota_{X_t}\omega_t\bigr)=0.
\]
This proves the first assertion.

In the equivariant case, first average \(\alpha\) over \(G\).
The normal splittings are invariant. Averaging \(f_t\) therefore
preserves its prescribed first jets, and makes \(X_t\) invariant.
The flow is equivariant. An invariant vector field is tangent to
the fixed locus, and its flow, starting at the identity, preserves
each fixed component. This also explains why no disjointness
condition between a fixed component and the \(S_j\) is needed.
\end{proof}

Only preservation as a set is asserted here. Equal pointwise
restrictions of \(\omega_0,\omega_1\) to a submanifold are unnecessary.
We will apply the lemma to divisors that are to be omitted, using
their complements after the isotopy.

\subsection{Projective degenerations}

\begin{lemma}[Transfer from smooth components]\label{geo:projective-transfer}
Let \(f:\mathcal X\to B\) be a flat projective family over a
smooth algebraic curve over \(\C\), with a marked point \(0\), and let
\(\mathcal L\) be a relatively ample rational polarization.
Let \(Y\) be a smooth projective irreducible component of
\(\mathcal X_0\), or a union of pairwise disjoint smooth projective
irreducible components.
The following statements hold after replacing \(B\)
by a Zariski-open neighborhood of \(0\); for symplectic transport
we subsequently work in a small complex analytic neighborhood.
\begin{enumerate}
\item For a sufficiently divisible and sufficiently large \(N\),
a full basis of
\(H^0(Y,\mathcal L^N|_Y)\) can be extended to sections defining a
relative projective embedding, with additional sections whose
restrictions to \(Y\) vanish. Consequently one can choose a
relative K\"ahler form whose restriction to \(Y\) is the
projective form from that full basis, divided by \(N\).
If a complex torus \((\C^*)^r\) acts holomorphically on \(Y\) and its action is
linearized on \(\mathcal L^N|_Y\), the basis may be chosen by
weights, making the restricted form invariant under the compact torus
\((S^1)^r\).
\item Let \(K\subset Y\) be compact and contained in the smooth
submersion locus of \(f\). For the chosen form there is, for all
sufficiently small nonzero \(t\), a symplectic embedding of a
neighborhood of \(K\) into the fiber \(\mathcal X_t\).
If \(K\) avoids a specified closed subset of \(\mathcal X\),
the embedding can be chosen with image avoiding that subset.
\item If a compact group acts on the family over \(B\) and the
total form is invariant, this transport is equivariant.
In particular, for a Hamiltonian circle action, it preserves
moment functions up to an additive constant on each connected
transport domain. If one transported point has the same prescribed
moment value in both fibers, this constant is zero.
\end{enumerate}
\end{lemma}

\begin{proof}
First clear denominators in the polarization. For large divisible
\(N\), relative Serre vanishing gives
\(R^1f_*(\mathcal I_Y\otimes\mathcal L^N)=0\);
see \cite[Tag~02O1]{Stacks}, or the relative ampleness and
vanishing theorem \cite[Theorem~1.7.6]{Lazarsfeld2004}.
Apply \(f_*\) to
\[
 0\longrightarrow\mathcal I_Y\otimes\mathcal L^N
 \longrightarrow\mathcal L^N
 \longrightarrow i_*(\mathcal L^N|_Y)\longrightarrow0,
\]
where \(i:Y\hookrightarrow\mathcal X\).
For a union of pairwise disjoint components, \(\mathcal I_Y\) is the
ideal of the whole union and the section space is the direct sum of the
component section spaces. The same lifting and local transport
arguments apply to that union.
The resulting surjection extends every section on \(Y\) near \(0\);
here \(f_*i_*(\mathcal L^N|_Y)\) is supported at \(0\), with fiber
\(H^0(Y,\mathcal L^N|_Y)\).
Take an affine neighborhood of \(0\), so these lifts may all be
represented by sections over that neighborhood.
Choose lifts \(s_0,\ldots,s_r\) of the desired basis. Choose also
sections \(t_0,\ldots,t_M\) of the same power which define a
relative closed immersion; eventual relative very ampleness over
the affine base follows from \cite[Tag~01VT]{Stacks} and properness.
Express each \(t_i|_Y\) in the chosen
basis, and subtract the corresponding linear combination of the
\(s_j\). The modified sections \(t_i'\) vanish on \(Y\).
The span of the \(s_j,t_i'\) contains the original \(t_i\), so this
enlarged system still defines a relative closed immersion.
On \(Y\), the resulting Fubini--Study potential is precisely
\(N^{-1}\log\sum_j|s_j|^2\) in a local frame.
For a linearized complex torus action, choose a basis of weight vectors.
Its compact torus characters preserve the standard diagonal Hermitian
norm, making this expression invariant under the compact torus.

The relative embedding into \(\PP^{r+M+1}\times B\) supplies a closed
real \((1,1)\)-form \(\Omega\), equal to the projective pullback
divided by \(N\), with an arbitrary positive base form added.
On the smooth locus it is K\"ahler, and its restriction to each
smooth fiber represents the designated polarization.
Adding the base form changes neither fiber forms nor the
horizontal lifts used below.

In the submersion locus define a lift of a base tangent vector by
requiring it to be \(\Omega\)-orthogonal to the fiber tangent space.
The fiber restriction of \(\Omega\) is symplectic, so the lift
exists uniquely. Lift a sufficiently short path from \(0\) to \(t\).
Compactness of \(K\) gives existence of its flow for a uniform
time, on a neighborhood of \(K\). Since \(\Omega\) is closed and
the contraction of the horizontal lift with \(\Omega\) vanishes
on fiber tangent vectors, this flow preserves the fiber forms.
It gives the claimed symplectic embedding. If \(K\) is disjoint
from a closed set, it has a neighborhood disjoint from that set;
shorten the transport uniformly so that its image remains in
this neighborhood.

For the last assertion, invariance of \(\Omega\) and uniqueness of
the horizontal lift imply equivariance. Write \(\Phi_t\) for the
transport and \(X\) for a circle generator. If
\(\iota_X\omega_t=-\dd\mu_t\), then
\[
 \dd(\mu_t\circ\Phi_t)
 =-\Phi_t^*(\iota_X\omega_t)
 =-\iota_X\omega_0
 =\dd\mu_0.
\]
The difference is constant on a connected domain, and a single
normalizing point determines it. When the circle acts on the
whole family, an invariant total form can be obtained by
averaging. If its restriction to \(Y\) was already invariant,
this averaging leaves that restriction unchanged.
\end{proof}

\begin{proposition}[Normal-cone transfer]\label{geo:transfer}
Let \(T\) be a smooth projective manifold, \(Z\subset T\) a smooth
closed submanifold of positive codimension, possibly disconnected,
and \(L\) an ample
rational polarization. Let
\[
 T'=\Bl_ZT,\qquad
 Y=\PP_Z(\mathbf1\oplus N_{Z/T}),\qquad
 D_\infty=\PP_Z(N_{Z/T})\subset Y.
\]
Write \(E_Z\) for the exceptional divisor on \(T'\); when \(Z\) is
a divisor, \(T'\simeq T\) and \(E_Z\) means \(Z\).
Suppose \(d>0\) is rational and
\[
 \pi^*L-dE_Z\quad\hbox{on }T',
 \qquad L|_Z+dU\quad\hbox{on }Y
\]
are ample.

Let \(F\subset T\) be a finite union of pairwise disjoint smooth
closed complex submanifolds. Assume the inverse image of
\(F\cap Z\) under \(Y\to Z\) is also a finite union of pairwise
disjoint smooth complex submanifolds; the empty union is allowed.
Let \(\omega_T\) be any K\"ahler form in class \(L\), and let
\(\omega_Y\) be a K\"ahler form in class \(L|_Z+dU\), invariant
under scalar rotation of the normal summand.
Then every compact subset of
\[
 Y\setminus\bigl(D_\infty\cup\pi_Y^{-1}(F\cap Z)\bigr)
\]
has a neighborhood admitting a symplectic embedding into
\((T\setminus F,\omega_T)\), with its given form \(\omega_Y\).
In particular any compact packing there transfers with
neighborhood embeddings.
\end{proposition}

\begin{proof}
Form the deformation to the normal cone
\cite[Section~5.1]{Fulton1998},
\[
 \mathcal X=\Bl_{Z\times0}(T\times\mathbb A^1)
 \longrightarrow\mathbb A^1.
\]
Its central fiber is the reduced union of \(T'\) and \(Y\).
The normal bundle of \(Z\times0\) in \(T\times\mathbb A^1\) is
\(N_{Z/T}\oplus\mathbf1\), giving the displayed description of
\(Y\). In the convention of lines,
\(\cO_{\mathcal X}(Y)|_Y=\cO_Y(-1)\).
Thus the rational polarization
\(\mathcal L=L-d[Y]\) restricts to the two classes in the
statement. The components meet along the exceptional divisor of
\(T'\), identified in \(Y\) with \(D_\infty\).
One can check the geometry in local coordinates
\((y,x_1,\ldots,x_r,t)\), with \(Z=\{x_1=\cdots=x_r=0\}\),
by blowing up the ideal \((x_1,\ldots,x_r,t)\).
In the \(t\)-chart, \(x_i=tv_i\); the exceptional component is
\(t=0\), and the family is a submersion there.
In the \(x_1\)-chart, write \(x_i=x_1v_i\) for \(i\ge2\)
and \(t=x_1s\). The central fiber is the reduced crossing
\(\{x_1s=0\}\), and its double locus is \(\{x_1=s=0\}\).
The other charts are identical. They show that the total space
is smooth, the family is flat, and the only nonsubmersion points
of the central fiber lie at this double locus.
The exceptional normal bundle and the restriction sign also
follow from the tautological line in the blowup construction;
see \cite[Tags~01OF, 02OS, and~0807]{Stacks}.

Ampleness on the two components implies ampleness on the central
fiber: the map from their disjoint union to that fiber is finite
and surjective, and ampleness descends under such maps.
Openness of relative ampleness then makes \(\mathcal L\) relatively
ample after shrinking about zero. Apply these statements to a
divisible integral multiple of \(\mathcal L\); precise forms are
\cite[Tags~0B5V and~0D2N]{Stacks}, and ampleness on components
and openness are also given by
\cite[Proposition~1.2.16 and Theorem~1.2.17]{Lazarsfeld2004}.
The nonzero fibers are canonically \(T\), with polarization \(L\).
The total space is smooth, and the family is a submersion along
\(Y\setminus D_\infty\).

The scalar action on \(Y\) linearizes on its polarization. By
Lemma~\ref{geo:projective-transfer}, choose a relative projective
form whose restriction \(\widehat\omega_Y\) is invariant.
Apply Lemma~\ref{geo:relative-moser} on \(Y\) to pass from
\(\omega_Y\) to \(\widehat\omega_Y\), preserving the indicated
vertical submanifolds. They are scalar-invariant. Moreover
\(D_\infty\) is a union of fixed components of scalar rotation and is
automatically preserved by the equivariant isotopy, including
at its intersections with the vertical submanifolds.

The image of the given compact set is now a compact subset of the
submersion locus, disjoint from the inverse image of \(F\) under
the natural map \(\mathcal X\to T\).
Lemma~\ref{geo:projective-transfer} transports a neighborhood into
a nearby \(T\setminus F\), with a K\"ahler form
\(\widehat\omega_T\) in class \(L\).
Finally apply Lemma~\ref{geo:relative-moser} on \(T\), preserving
the components of \(F\), to identify \(\widehat\omega_T\) with
\(\omega_T\). Composing these maps proves the proposition.
All maps are defined on neighborhoods of the relevant compact
sets; these neighborhoods can be shrunk at each composition.
\end{proof}

\subsection{An explicit normal-bundle form}

\begin{lemma}[Equal positive normal summands]\label{geo:bundle}
Let \(A\) be an ample line bundle on a smooth projective manifold
\(Z\), equipped with a Hermitian metric of Chern form
\(\sigma>0\). Suppose
\[
 N=A^b\oplus\cdots\oplus A^b
\]
has \(\ell\) summands, where \(b\) is a positive integer.
If \(d>0\) and \(bd<1\), the class
\(\sigma+dU\) on \(Y=\PP_Z(\mathbf1\oplus N)\) has an invariant
K\"ahler representative with the following affine description:
\begin{equation}\label{geo:bundle-form}
 \omega_Y=\sigma+\dd\left(\sum_{j=1}^{\ell}p_j\alpha_j\right),
 \qquad p_j\ge0,\quad \sum_jp_j<d.
\end{equation}
Here \(\alpha_j\) is the angular connection of period one on the
\(j\)-th normal line, with curvature \(-b\sigma\), and the fiber
coordinates have standard moments \(p_j\).
The expression \(p_j\alpha_j\) is smooth across its zero axis.
For rational \(d\), this also proves ampleness of the indicated
rational polarization.

If \(Z=C\) is a curve, then over a compact bordered subsurface
\(\Sigma\) the model, truncated on any compact moment polytope
strictly inside \(\sum p_j<d\), admits a toric horizontal primitive
as in Lemma~\ref{surf:packing}, with
\[
 \kappa(p)=\left(1-b\sum_jp_j\right)\sigma.
\]
For bordered exhaustions of \(C\) minus finitely many points, the
integrated areas converge uniformly on such a polytope to
\[
 H(p)=(\deg A)\left(1-b\sum_jp_j\right).
\]
\end{lemma}

\begin{proof}
On the affine chart \(N\subset\PP(\mathbf1\oplus N)\), let
\(w_1,\ldots,w_\ell\) be the normal coordinates, with their induced
Hermitian norms. Set
\[
 \omega_Y=\sigma+
 \dd\dd^c\left[d\log\left(1+\sum_j\lVert w_j\rVert^2\right)\right].
\]
The expression in brackets is the local weight of the dual
tautological metric, multiplied by \(d\), so its curvature
extends globally and represents \(dU\).
The moments of the normal rotations are
\begin{equation}\label{geo:bundle-moments}
 p_j=\frac{d\lVert w_j\rVert^2}
 {1+\sum_l\lVert w_l\rVert^2}.
\end{equation}
Keeping the phases and using radii \(\sqrt{p_j/\pi}\) identifies
each affine fiber with the standard open ball
\(\sum_jp_j<d\). The inverse radial change is smooth away from
the upper boundary, including all zero coordinates.
Writing the metric connection in these coordinates yields
\eqref{geo:bundle-form}. Equivalently, differentiating the
potential gives \(\sum_jp_j\alpha_j\). In a unitary local frame,
\(\alpha_j=\dd\theta_j+A_j\), and
\(p_j\dd\theta_j\) is the standard smooth Liouville form of the
complex coordinate, while \(p_jA_j\) is plainly smooth.

At any base point choose a holomorphic Chern frame whose metric
has zero first derivative at that point. The form then splits
into a positive vertical Fubini--Study form and the horizontal
form
\[
 \left(1-b\sum_jp_j\right)\sigma.
\]
This remains valid on the projective infinity chart, where
\(\sum_jp_j=d\). The horizontal factor is bounded below by
\(1-bd>0\), so the extended form is K\"ahler.
For rational \(d\), a divisible integral multiple is the
curvature of a positive Hermitian line bundle, which is ample
by the Kodaira embedding theorem
\cite[\S1.2.C, p.~43]{Lazarsfeld2004}.

For the curve assertion, every complex line bundle on a bordered
surface is smoothly trivial. Choose unitary frames on a slightly
larger bordered neighborhood of \(\Sigma\), so
\(\alpha_j=\dd\theta_j+A_j\) there. Choose a base primitive
\(\gamma\) of \(\sigma\). Then
\[
 \omega_Y=\sum_j\dd p_j\wedge\dd\theta_j+
 \dd\Gamma_0,\qquad
 \Gamma_0=\gamma+\sum_jp_jA_j,
\]
where \(\Gamma_0\) is horizontal and toric. Its base differential
is \((1-b\sum_jp_j)\sigma\).
The assertion about integrated areas follows from
\(\int_\Sigma\sigma\to\int_C\sigma=\deg A\).
The moment coefficient is continuous and bounded on the fixed
compact polytope, so convergence is uniform.
\end{proof}

\subsection{The target ball and its exterior shell}

The complement of a hyperplane \(H_\infty\) in unit projective
space is the open capacity-one ball. More explicitly, in affine
coordinates \(w\in\C^n\) the Fubini--Study moments and phases are
\[
 p_j=\frac{|w_j|^2}{1+\sum_l|w_l|^2},\qquad
 \theta_j=\frac{\arg w_j}{2\pi}.
\]
The form is \(\sum_j\dd p_j\wedge\dd\theta_j\), and the map
\begin{equation}\label{geo:affine-ball}
 w\longmapsto
 \frac{w}{\sqrt{\pi(1+\sum_l|w_l|^2)}}
\end{equation}
is a symplectomorphism onto
\(\{z:\pi\sum_j|z_j|^2<1\}\). The identity of forms is first
verified where all coordinates are nonzero and then extends
smoothly across the axes.

\begin{lemma}[A blowup complement is a shell]\label{geo:shell}
Let \(n\ge2\), let \(o\in\PP^n\setminus H_\infty\), and write
\[
 T=\Bl_o\PP^n,\qquad H=\pi^*c_1(\cO_{\PP^n}(1)).
\]
Let \(E\) be the exceptional divisor and use \(H_\infty\) also
for its proper transform. For \(0<a<1\), there is a K\"ahler
form \(\omega_a\) in class \(H-aE\) for which
\[
 (T\setminus(E\cup H_\infty),\omega_a)
 \simeq
 \left(\left\{z\in\C^n:
 a<\pi\sum_j|z_j|^2<1\right\},\omega_0\right).
\]
The same symplectomorphism type holds for every K\"ahler form
in this class. In addition, \(AH-BE\) is ample as a rational
polarization whenever \(A>B>0\) are rational.
\end{lemma}

\begin{proof}
Choose homogeneous coordinates with
\(o=[1:0:\cdots:0]\) and \(H_\infty=\{z_0=0\}\).
On \(\PP^n\times\C\) let the circle act by
\[
 ([z_0:z'],v)\longmapsto
 ([z_0:e^{2\pi it}z'],e^{-2\pi it}v).
\]
For the product of unit Fubini--Study and the standard form on
\(\C\), its Hamiltonian is
\[
 F=\mu-\pi|v|^2,\qquad
 \mu=\frac{|z'|^2}{|z_0|^2+|z'|^2}.
\]
The level \(F=a\) is compact, regular, and the circle acts freely
there. Indeed, if \(v\ne0\) the action is free on that coordinate;
if \(v=0\), then \(\mu=a\in(0,1)\), where the projective scalar
action is free. Reduction therefore gives a smooth compact
K\"ahler manifold. This is Lerman's symplectic cut
\cite[Section~1.1]{Lerman1995}; its K\"ahler interpretation is
described in \cite[Section~2]{BurnsGuilleminLerman2002}.

For completeness, its complex model can be identified directly.
On the stable set
\[
 z'\ne0,\qquad (z_0,v)\ne(0,0),
\]
the holomorphic map
\[
 ([z_0:z'],v)\longmapsto
 \bigl([z_0:vz'],[z']\bigr)
\]
takes values in the graph resolution of the projection from
\(o\), which is \(\Bl_o\PP^n\), and its fibers are the free
complex scalar orbits. On each such orbit write the positive
real scalar as \(\lambda\). The function
\[
 \frac{\lambda^2|z'|^2}{|z_0|^2+\lambda^2|z'|^2}
 -\frac{\pi|v|^2}{\lambda^2}
\]
is strictly increasing from a value at most zero, or from
\(-\infty\), to the limit one. Consequently it takes the value
\(a\) exactly once. Each complex orbit meets the level in one
circle, and the K\"ahler quotient has the indicated complex
blowup structure.

On the retained open region \(\mu>a\), choose the representative
\[
 v=\sqrt{(\mu-a)/\pi}>0.
\]
The auxiliary complex-line form pulls back to zero on this
real slice, so the original Fubini--Study form is unchanged.
At \(v=0\), reduction of the level \(\mu=a\) gives \(E\simeq
\PP^{n-1}\) with line area \(a\), as is seen from its moment
simplex \(\sum_jp_j=a\). The proper transform of \(H_\infty\)
is retained in the open region and has unit line area.
These two restrictions determine the class \(H-aE\):
\(H|_E=0\) and \(E|_E=-c_1(\cO_E(1))\).
Removing \(E\) and \(H_\infty\), and then applying
\eqref{geo:affine-ball}, gives precisely the stated open shell.

For any other K\"ahler representative of \(H-aE\),
Lemma~\ref{geo:relative-moser} gives an identification preserving
the disjoint complex divisors \(E,H_\infty\), and therefore their
complement.

Finally the graph model is a closed subvariety of
\(\PP^n\times\PP^{n-1}\). The pullbacks of the two hyperplane
classes are \(H\) and \(H-E\), respectively, since projection
from \(o\) is defined by hyperplanes through \(o\).
The restriction of a positive rational product polarization is
ample. The identity
\[
 AH-BE=(A-B)H+B(H-E)
\]
therefore proves the last assertion.
\end{proof}
